% ===================================================================
%  TABELL Nr 15 — Marknaden. — Matvarorna.
%  Traduction 2026 d'apres texto/io, controlee sur texto/fr et
%  texto/nl. Consigne : voir texto/sv/10-tabell-01.tex.
%
%  UN TABLEAU DE LISTES, ou le renvoi suit le nom de la denree et ou
%  la phrase n'a presque pas d'occasion de se retourner. La seule
%  precaution est celle du tableau 14 : partout ou l'ido postpose son
%  possesseur et ou le possede porte un numero, on ecrit « troskeln
%  till bageriet (78) » et non « bageriets troskel », pour que (80) et
%  (78) restent dans l'ordre.
% ===================================================================

%%K t15-apar-1 apar
\VUsaut{4.0mm}
\VUtitre{15.0pt}{60}{TABELL Nr 15}
\VUsaut{3.0mm}

%%K t15-tit-1 sub
\VUpk{11.4pt}{40}{Marknaden. --- Matvarorna.}
\VUsaut{3.0mm}

%%K t15-01-1 p
1. --- De flesta av de handlande som levererar oss de varor som är nödvändiga
för vår föda samlas under \VUgras{hallen} \textsuperscript{(1)} på den
\VUgras{täckta marknaden} eller på de angränsande gatorna.

%%K t15-02-1 p
2. --- \VUgras{Charkuteristen} \textsuperscript{(2)}, som säljer \VUgras{skinka}
\textsuperscript{(3)}, \VUgras{blodkorv} \textsuperscript{(4)},
\VUgras{korv} \textsuperscript{(5)}, \VUgras{ostkorv} \textsuperscript{(6)} och
\VUgras{fläsk} \textsuperscript{(7)}, står bredvid \VUgras{smör- och
ostförsäljerskan} \textsuperscript{(8)}, vars skyltning är försedd med
\VUgras{holländska ostar} \textsuperscript{(9)} med röd skorpa, med
\VUgras{camembertostar} \textsuperscript{(10)}, med \VUgras{schweizerostbitar}
\textsuperscript{(11)} och med färska \VUgras{smörklickar}
\textsuperscript{(12)}.

%%K t15-03-1 p
3. --- Framför henne bjuder en \VUgras{grönsaksförsäljerska}
\textsuperscript{(13)} sina kunder vackra \VUgras{tomater}
\textsuperscript{(14)}, \VUgras{blomkål} \textsuperscript{(15)},
\VUgras{sallad} \textsuperscript{(16)}, \VUgras{huvudkål}
\textsuperscript{(17)}, \VUgras{squash} \textsuperscript{(19)},
\VUgras{meloner} \textsuperscript{(18)} och till och med \VUgras{svampar}
\textsuperscript{(20)}. Men en \VUgras{ölutkörare}, som har stannat sin
\VUgras{fraktvagn} \textsuperscript{(21)} lastad med \VUgras{ölfat}
\textsuperscript{(22)} just framför hennes skyltning, hindrar hennes kunder från
att komma nära. En trädgårdsmästarhustru bär till henne en stor korg
\VUgras{kronärtskockor} \textsuperscript{(23)}.

%%K t15-tit-2 sub
\VUsaut{4.0mm}
\VUpk{10.2pt}{40}{Att sälja och att köpa.}
\VUsaut{3.0mm}

%%K t15-04-1 p
4. --- På marknaden är trängseln stor, ty timmen för \VUgras{auktionen}
\textsuperscript{(24)} har kommit. \VUgras{Husmödrarna} \textsuperscript{(26)},
som kommer dit för sina inköp, stannar i förbigående framför
\VUgras{komagsförsäljerskans} \textsuperscript{(25)} stånd, för att köpa
\VUgras{komage} eller ett \VUgras{kalvhuvud}.

%%K t15-05-1 p
5. --- En tjock \VUgras{kock} \textsuperscript{(27)} prutar på en mycket vacker
\VUgras{tonfisk} hos \VUgras{fiskförsäljerskan} \textsuperscript{(28)} som
höjer sina (försäljnings)priser efter behag. Hon har fått en stor mängd
\VUgras{ålar, sjötungor, foreller} och \VUgras{karpar}. Hennes \VUgras{torsk}
och hennes \VUgras{musslor} är mycket färska.

%%K t15-tit-3 sub
\VUsaut{4.0mm}
\VUpk{10.2pt}{40}{Billiga och dyra varor.}
\VUsaut{3.0mm}

%%K t15-06-1 p
6. --- \VUgras{Fågelhandlerskan} \textsuperscript{(29)}, som är installerad
bredvid henne, är mycket samvetsgrann. När hon säljer en \VUgras{kalkon, en
gås, en anka} eller en vanlig \VUgras{höna}, begär hon bara det rätta priset.

%%K t15-07-1 p
7. --- I hörnet av torget, på första våningen, har sedan kort tid en
\VUgras{linneförsäljerska} \textsuperscript{(30)} slagit sig ner hos vilken
köperskorna alltid är säkra på att finna ett mycket vackert förråd av
\VUgras{bordsdukar, servetter, förkläden} och \VUgras{handdukar}. På samma
våning har en \VUgras{knivsmed} \textsuperscript{(31)} inrättat sin verkstad.
Han slipar försäljerskornas \VUgras{knivar} i hallen och gör åt
trädgårdsmästarna \VUgras{trädgårdsknivar} av det hårdaste \VUgras{stål}.

%%K t15-tit-4 sub
\VUsaut{4.0mm}
\VUpk{10.2pt}{40}{Specerihandeln.}
\VUsaut{3.0mm}

%%K t15-08-1 p
8. --- Under hans verkstad har sedan kort tid en stor livsmedelsaffär öppnats.
Det är redan inte längre den enkla och obetydliga \VUgras{specerihandeln}
\textsuperscript{(32)} från förr, som bara sålde de gångbara varorna,
\VUgras{te} \textsuperscript{(33)}, \VUgras{choklad} \textsuperscript{(34)},
\VUgras{sockertopp} \textsuperscript{(37)} och \VUgras{tvål}
\textsuperscript{(38)}. Där säljer man varje vara, \VUgras{knapriga skorpor},
\VUgras{konserver} \textsuperscript{(35)}, \VUgras{likörer}
\textsuperscript{(36)}, \VUgras{rom, konjak, kirsch, anisett} och
\VUgras{curaçao}. Där finner man vilt : \VUgras{hare} \textsuperscript{(39)},
\VUgras{trastar} \textsuperscript{(40)}, \VUgras{rapphöns}
\textsuperscript{(41)}, \VUgras{vaktlar} \textsuperscript{(42)}, \VUgras{vildand}
\textsuperscript{(43)} eller \VUgras{fasan} \textsuperscript{(44)}, lika lätt
som utländska frukter såsom \VUgras{bananer} \textsuperscript{(46)} och
\VUgras{ananaser}.

%%K t15-09-1 p
9. --- En mycket tjänstvillig \VUgras{specerihandelsgesäll}
\textsuperscript{(45)}, står redo att ge vad man önskar : \VUgras{sylt}
\textsuperscript{(47)} av vinbär eller av kvitten, \VUgras{ister}
\textsuperscript{(48)}, \VUgras{ättiksgurkor} \textsuperscript{(49)} eller
saltade \VUgras{sardiner} \textsuperscript{(50)}.

%%K t15-09-2 p
Det är mycket bekvämt för \VUgras{kunderna} \textsuperscript{(51)}, som, vid
sidan av de mest gångbara varorna, finner de sällsyntaste förstlingarna och de
smakligaste frukterna : saftiga \VUgras{päron} \textsuperscript{(52)},
\VUgras{plommon} \textsuperscript{(53)}, \VUgras{vindruvor}
\textsuperscript{(54)}, \VUgras{persikor} \textsuperscript{(55)},
\VUgras{mandlar} \textsuperscript{(56)} och \VUgras{nötter}
\textsuperscript{(57)}.

%%K t15-tit-5 sub
\VUsaut{4.0mm}
\VUpk{10.2pt}{40}{Kundkretsen.}
\VUsaut{3.0mm}

%%K t15-10-1 p
10. --- \VUgras{Riset} \textsuperscript{(58)} och \VUgras{linserna}
\textsuperscript{(59)} står bredvid lådan med rökta \VUgras{sillar}
\textsuperscript{(60)}; \VUgras{potatisen} \textsuperscript{(61)} och
\VUgras{bönan} \textsuperscript{(63)} står bredvid \VUgras{äpplena}
\textsuperscript{(62)}, \VUgras{fikonen} \textsuperscript{(64)},
\VUgras{katrinplommonen} \textsuperscript{(65)} och \VUgras{hasselnötterna}
\textsuperscript{(66)}. Man finner i den boden färska \VUgras{ägg}
\textsuperscript{(67)} och \VUgras{oliver} \textsuperscript{(68)}; därför fylls
\VUgras{korgarna} \textsuperscript{(70)} och \VUgras{nätkassarna}
\textsuperscript{(73)} med rasande fart.

%%K t15-11-1 p
11. --- \VUgras{Kaffet} \textsuperscript{(72)} från den utmärkta firman har bland
\VUgras{tjänarinnorna} \textsuperscript{(69)} och \VUgras{kokerskorna} vunnit ett
välförtjänt rykte, ty den gamle \VUgras{specerihandlaren} \textsuperscript{(71)}
rostar det själv med den största omsorg.

%%K t15-tit-6 sub
\VUsaut{4.0mm}
\VUpk{10.2pt}{40}{Skådespelen.}
\VUsaut{3.0mm}

%%K t15-12-1 p
12. --- Alla går inte till marknaden för att förse sig. I den pittoreska
trafiken på den lilla gata som leder till hallen ser man de mest olikartade
människor. Bredvid en vänlig \VUgras{slaktargesäll} \textsuperscript{(75)}, som
skjuter en \VUgras{handkärra} \textsuperscript{(74)} framför sig, se här två
permitterade soldater som är mycket stolta över att visa sin lysande uniform :
\VUgras{husaren} \textsuperscript{(76)} med en blå \VUgras{dolma} och en
\VUgras{schakå med plym}, och \VUgras{zuavkorpralen}, med vida byxor och
huvudet täckt av en \VUgras{fez}.

%%K t15-13-1 p
13. --- \VUgras{Bagaren} \textsuperscript{(80)}, som står på tröskeln till
\VUgras{bageriet} \textsuperscript{(78)} har just knådat degen till sitt
\VUgras{bröd} \textsuperscript{(79)} och tagit ut ur ugnen \VUgras{giffarna}
\textsuperscript{(81)}, \VUgras{småfranskorna} \textsuperscript{(82)} och de
\VUgras{runda bröden} \textsuperscript{(83)} som ligger i skyltfönstret.

%%K t15-14-1 p
14. --- Ovanför bageriet bor en skicklig \VUgras{linnesömmerska}
\textsuperscript{(84)} som har flera \VUgras{broderskor} \textsuperscript{(85)}
i sin tjänst. Bredvid henne bor en \VUgras{tvätterska-stryklerska}
\textsuperscript{(86)}, som stärker \VUgras{linnet} \textsuperscript{(88)}
sedan hon har tvättat och torkat det, och som stryker det med ett
\VUgras{strykjärn} \textsuperscript{(87)}.

%%K t15-tit-7 sub
\VUsaut{4.0mm}
\VUpk{10.2pt}{40}{Kött, fisk och grönsaker.}
\VUsaut{3.0mm}

%%K t15-15-1 p
15. --- I \VUgras{slakteriet} \textsuperscript{(89)}, som ligger på
bottenvåningen, säljer man alla slags \VUgras{kött, fårkött}
\textsuperscript{(90)}, \VUgras{oxkött} \textsuperscript{(94)} eller
\VUgras{kalvkött} \textsuperscript{(92)} såsom \VUgras{fårstekar, biffstek,
stek} och \VUgras{kotletter}. \VUgras{Slaktargesällen} \textsuperscript{(93)}
skär allt det köttet och lägger \VUgras{märgbenen} \textsuperscript{(96)} för
sig. Vid slakteriets dörr står en \VUgras{ostronförsäljerska}
\textsuperscript{(97)} som säljer \VUgras{ostron} \textsuperscript{(98)}. Hon
öppnar dem om man önskar det.

%%K t15-16-1 p
16. --- Alla de handlande betalar en näringsskatt eller en ståndsavgift; därför
förföljer \VUgras{poliskonstapeln} \textsuperscript{(100)} utan medlidande de
stackars \VUgras{grönsaksmånglerskorna} \textsuperscript{(99)} som konkurrerar
med dem genom att med sina kärror gå omkring och sälja \VUgras{morötter,
rädisor, rovor, purjolök} eller \VUgras{sparrisknippen, färska ärter, spenat,
syra, vattenkrasse} och \VUgras{persilja}.

%%K t15-tit-8 sub
\VUsaut{4.0mm}
\VUpk{10.2pt}{40}{Akta dig för poliskonstapeln!}
\VUsaut{3.0mm}

%%K t15-17-1 p
17. --- Pojken som säljer \VUgras{vitlök} \textsuperscript{(101)} och
\VUgras{lök}, vet mycket väl, att också han inte får sälja sina varor vid
marknaden.

%%K t15-17-1 p suite
Han flyr, och riskerar därvid att välta korgen med \VUgras{krabbor},
\VUgras{kräftor} och \VUgras{räkor} hos \VUgras{försäljaren}
\textsuperscript{(102)} av \VUgras{languster} och \VUgras{humrar}.

%%K t15-18-1 p
18. --- Alla rör sig och gör mycket buller; men det hindrar inte att man tydligt
hör den kringvandrande \VUgras{glasmästarens} \textsuperscript{(103)} röst,
eller ropet från \VUgras{kolhandlaren} \textsuperscript{(104)} som just några
ögonblick har lämnat sin kärra för att leverera en \VUgras{säck kol}
\textsuperscript{(105)}.

\VUsaut{6.0mm}
\VUfilet{22mm}
